%%%%%%%%%%%%%%%%%%%%%%%%%%%%%%%%%%%%%%%%%%%%%%%%%%%%%%%%%%%%%%%%%%%%%%%%%%%%%%%%
%                   CHAPTER 6: CONCLUSIONS & FUTURE WORK                       %
%%%%%%%%%%%%%%%%%%%%%%%%%%%%%%%%%%%%%%%%%%%%%%%%%%%%%%%%%%%%%%%%%%%%%%%%%%%%%%%%

\chapter{Conclusions \& Future Work}\label{section:conclusions}

This dissertation addressed the architectural tension between data privacy rights and mandatory healthcare record preservation within distributed Electronic Health Record (EHR) environments. By decoupling physical data persistence from cryptographic accessibility, CryptoShred Health demonstrates that permanent erasure under GDPR Article~17 can be enforced across distributed databases, caching tiers, event brokers, and immutable WORM backups without requiring destructive mutations of append-only storage blocks.

%%%%%%%%%%%%%%%%%%%%%%%%%%%%%%%%%%%%%%%%%%%%%%%%%%%%%%%%%%%%%%%%%%%%%%%%%%%%%%%%
%            SECTION 6.1: SUMMARY OF CONTRIBUTIONS & ACHIEVEMENTS              %
%%%%%%%%%%%%%%%%%%%%%%%%%%%%%%%%%%%%%%%%%%%%%%%%%%%%%%%%%%%%%%%%%%%%%%%%%%%%%%%%

\section{Summary of Contributions \& Research Achievements}\label{section:contributions}

The primary research objective was to design, implement, and empirically validate an end-to-end zero-plaintext EHR platform that achieves constant-time ($\mathcal{O}(1)$) cryptographic erasure, verifiable post-quantum auditability, and high-throughput query performance. By systematically executing the \textbf{Question $\rightarrow$ Experiment $\rightarrow$ Results} methodology established in Chapter~\ref{section:introduction}, this research resolved all four core Research Questions ($RQ_1$--$RQ_4$).

Table~\ref{tab:rq-resolution-matrix} provides a synthesis of the research questions, governing hypotheses, experimental testbeds, and observed findings.

\begin{table}[H]
  \centering
  \caption{Research Questions Resolution and Empirical Validation Matrix.}
  \label{tab:rq-resolution-matrix}
  \begin{tabular}{p{3.2cm} p{4.2cm} p{5.6cm} c}
    \toprule
    \textbf{Research Question} & \textbf{Target Hypothesis ($H$)} & \textbf{Empirical Finding (Measured)} & \textbf{Verdict} \\
    \midrule
    \textbf{$RQ_1$: Erasure Scalability} & 
    $H_1$: Constant-time $\mathcal{O}(1)$ shredding vs $\mathcal{O}(N)$ SQL cascade; sub-$10\text{~ms}$ latency & 
    End-to-end crypto-shredding completes in $5.63\text{--}7.65\text{~ms}$ ($4\text{--}5\times$ faster than PostgreSQL cascades at $25.19\text{--}39.12\text{~ms}$); microbenchmarks confirm $\mathcal{O}(1)$ invariance ($59.50\,\mu\text{s}$) & 
    \textbf{Validated} \\
    \addlinespace[0.6em]
    \textbf{$RQ_2$: Post-Quantum Proofs} & 
    $H_2$: Sub-$5\text{~ms}$ Merkle proof verification under ML-DSA-65 & 
    Merkle leaf insertion in $0.13\,\mu\text{s}$; full proof verification in $2.48\text{~ms}$ across $100{,}000$ tree leaves & 
    \textbf{Validated} \\
    \addlinespace[0.6em]
    \textbf{$RQ_3$: Snapshot Paradox} & 
    $H_3$: Linear-time purge ($\mathcal{O}(k)$), sub-$5\text{~ms/key}$, $<50\text{~ms}$ for drift cohorts ($k \le 30$) & 
    Tombstone reconciler purged all $25/25$ restored keys in $42.3\text{~ms}$ ($1.69\text{~ms/key}$) during boot audit & 
    \textbf{Validated} \\
    \addlinespace[0.6em]
    \textbf{$RQ_4$: Concurrency \& Leakage} & 
    $H_4$: $>50{,}000\text{~req/s}$ reads with $0.000\%$ post-shred leakage & 
    Sustained $96{,}820.5\text{~req/s}$ on cached reads with $0.000\%$ plaintext leakage across $2{,}377{,}883$ adversarial queries & 
    \textbf{Validated} \\
    \bottomrule
  \end{tabular}
\end{table}

The empirical findings of this research yield four foundational conclusions:
\begin{enumerate}
    \item \textbf{Constant-Time $\mathcal{O}(1)$ Cryptographic Erasure:} In production deployments, destroying a 256-bit Key Encryption Key (KEK) inside HashiCorp Vault enforces permanent data erasure end-to-end across all distributed tiers in $5.63\text{--}7.65\text{~ms}$. This achieves an honest $4\text{--}5\times$ system speedup over PostgreSQL transactional cascading deletions ($25.19\text{--}39.12\text{~ms}$) for patient histories of $N = 10{,}000$ records, while isolated CPU microbenchmarks confirm invariant $\mathcal{O}(1)$ computational scaling ($59.50\,\mu\text{s}$). This guarantees compliance with GDPR Article~17 without modifying immutable physical storage blocks.

    \item \textbf{Dynamic Retention Adaptation:} By calculating legal retention deadlines strictly over active, non-shredded clinical encounters ($\mathcal{V}_{\text{active}}(p)$), erasing an individual visit dynamically updates the patient's overall retention horizon. The governance engine records legal justification codes directly within signed erasure certificates, enabling automated compliance auditing without manual human intervention.

    \item \textbf{Automated Disaster Recovery and Resurrection Defense:} The dual-source tombstone reconciler resolves the KMS snapshot restoration paradox. By cross-referencing restored Vault key rings against both database state and immutable external WORM deletion receipts, the reconciler purged 100\% ($25/25$) of resurrected master keys in $42.3\text{~ms}$ ($1.69\text{~ms}$ per key) during boot audits, prior to opening public network listeners.

    \item \textbf{Post-Quantum Verifiability under High Concurrency:} Dual-signing erasure certificates with classical RSA-2048 and post-quantum ML-DSA-65 within an append-only binary Merkle tree enables independent third-party auditors to verify inclusion proofs in $2.48\text{~ms}$ across $100{,}000$ historical records ($\mathcal{O}(\log N)$). Under peak clinical concurrency (500 virtual users), the architecture sustains $96{,}820.5\text{~requests/s}$ on cached reads while maintaining a strict $0.000\%$ plaintext leakage rate across $2{,}377{,}883$ read requests targeting shredded entities.
\end{enumerate}

%%%%%%%%%%%%%%%%%%%%%%%%%%%%%%%%%%%%%%%%%%%%%%%%%%%%%%%%%%%%%%%%%%%%%%%%%%%%%%%%
%               SECTION 6.2: LIMITATIONS & TECHNICAL TRADE-OFFS                %
%%%%%%%%%%%%%%%%%%%%%%%%%%%%%%%%%%%%%%%%%%%%%%%%%%%%%%%%%%%%%%%%%%%%%%%%%%%%%%%%

\section{Limitations \& Technical Trade-Offs}\label{section:limitations}

While the architecture demonstrates measurable performance and verifiable privacy guarantees, several operational trade-offs and technical limitations warrant critical examination:

\begin{itemize}
    \item \textbf{Trusted Computing Base (TCB) Scope and Host Compromise:} The application server resides within the platform's Trusted Computing Base. While cryptographic erasure guarantees that destroyed keys and their associated ciphertexts remain computationally irrecoverable even against an adversary with host-level root privileges, active (non-shredded) patient data remains vulnerable to memory inspection if the application host is fully compromised. Protecting active keys in memory against root adversaries requires hardware-isolated enclaves.

    \item \textbf{Key Persistence in KMS Disaster Recovery Snapshots:} Backups of the HashiCorp Vault Raft cluster contain encrypted key rings protected by Vault barrier encryption. Restoring an outdated KMS snapshot can re-introduce previously destroyed keys into active storage. While our WORM tombstone reconciler detects and purges these resurrected keys upon startup, production deployments must restrict KMS snapshot access and enforce Shamir secret sharing thresholds ($k \ge 3$) for unsealing operations.

    \item \textbf{Volatile L2 Cache Semantics:} To uphold the zero-plaintext architecture, Redis operates strictly as an in-memory L2 cache within the volatile RAM tier of the TCB. Disk persistence is disabled (\texttt{save ""}, \texttt{appendonly no}), entries are bounded by a 15-minute TTL, and the crypto-shredding service proactively invalidates cached patient objects upon key destruction. Exposure to cold storage persistence is thereby eliminated, though residual plaintext exists in volatile RAM during active caching windows.

    \item \textbf{Internal Network Transport and Unseal Configuration:} In the prototype evaluation testbed, communication between Spring Boot and Vault traversed an internal Docker bridge network without TLS, and Vault operated with a single unseal key share for automation. Production deployments require mutual TLS (mTLS) across all service boundaries, AppRole authentication with tightly scoped token policies, and multi-party Shamir threshold configurations.

    \item \textbf{Uniform Statutory Retention Horizon:} The prototype retention engine evaluates retention periods using a single global parameter ($\Delta_{\text{retention}}$). In clinical practice, statutory requirements vary by clinical specialty, patient age, and jurisdiction. For example, HIPAA 45 CFR \S164.316 mandates 6-year retention specifically for privacy and compliance documentation, whereas clinical medical records are governed by state-level statutes and hospital policies.

    \item \textbf{JVM Memory Zeroization Guarantees:} While the backend explicitly clears primitive byte arrays via \texttt{Arrays.fill(dekBytes, (byte) 0)} within mandatory \texttt{finally} blocks, the managed Java Virtual Machine controls heap allocations without direct pointer destruction. Intermediate \texttt{String} and JSON object instances created during serialization persist in heap memory until reclaimed by the garbage collector. Furthermore, standard HotSpot runtimes do not execute \texttt{mlockall()}; anti-swap protection relies on host container configuration (\texttt{--memory-swap = --memory}).

    \item \textbf{Blind Index Information Leakage via Frequency Analysis:} Deterministic HMAC-SHA256 blind indexing enables constant-time exact-match lookups but cannot support range or fuzzy search predicates. Because deterministic hashing maps identical inputs to identical digest tokens, an adversary obtaining a full database dump could attempt frequency analysis against low-entropy demographic fields (such as common surnames within small regional clinics).

    \item \textbf{Post-Quantum Signature Storage Expansion:} While ML-DSA-65 signatures ensure long-term resistance against quantum cryptanalysis, lattice-based cryptography introduces signature payload expansion. An ML-DSA-65 signature requires $3{,}309\text{~bytes}$, compared to $256\text{~bytes}$ for classical RSA-2048, resulting in an approximate $13\times$ storage expansion for signed Merkle certificates.
\end{itemize}

%%%%%%%%%%%%%%%%%%%%%%%%%%%%%%%%%%%%%%%%%%%%%%%%%%%%%%%%%%%%%%%%%%%%%%%%%%%%%%%%
%                         SECTION 6.3: FUTURE WORK                             %
%%%%%%%%%%%%%%%%%%%%%%%%%%%%%%%%%%%%%%%%%%%%%%%%%%%%%%%%%%%%%%%%%%%%%%%%%%%%%%%%

\section{Future Work}\label{section:future-work}

Several promising research directions emerge from this architecture:

\begin{enumerate}
    \item \textbf{External Ledger Anchoring for Merkle Roots:} Periodically anchoring Merkle root digests into external immutable ledgers (such as RFC 6962 Certificate Transparency logs or public distributed ledgers) would provide verifiable chronological timestamping beyond local infrastructure boundaries.

    \item \textbf{Granular Multi-Jurisdiction Retention Engines:} Extending the retention scheduler into a declarative policy engine supporting department-level, diagnosis-specific, and pediatric retention policies across international regulatory frameworks (UK NHS Code, HIPAA, and Romanian Law 95/2006).

    \item \textbf{Large-Scale Synthetic Clinical Evaluation:} Evaluating platform ingestion, blind indexing, and crypto-shredding pipelines against comprehensive clinical datasets generated by the open-source HL7 FHIR Synthea engine, modeling complex multi-decade patient histories across thousands of encounters.

    \item \textbf{Hardware-Isolated Secure Enclaves (Confidential Computing):} Moving DEK unwrapping, payload decryption, and cryptographic zeroization into hardware enclaves (such as AWS Nitro Enclaves or Intel SGX) would protect active keys from memory scraping, isolating cryptographic operations from host operating system administrators.

    \item \textbf{Medical Analytics Over Encrypted Clinical Data:} Integrating functional encryption or homomorphic primitives to enable epidemiological aggregation and clinical analytics directly across encrypted records without requiring intermediate plaintext reconstruction in application memory.
\end{enumerate}

%%%%%%%%%%%%%%%%%%%%%%%%%%%%%%%%%%%%%%%%%%%%%%%%%%%%%%%%%%%%%%%%%%%%%%%%%%%%%%%%
%                        SECTION 6.4: FINAL REMARKS                            %
%%%%%%%%%%%%%%%%%%%%%%%%%%%%%%%%%%%%%%%%%%%%%%%%%%%%%%%%%%%%%%%%%%%%%%%%%%%%%%%%

\section{Final Remarks}\label{section:final-remarks}

The tension between a patient's fundamental right to erasure and the legal obligation to maintain medical history stems from the physical deletion paradigm of traditional database systems. Attempting to enforce data erasure by overwriting physical storage blocks inevitably compromises distributed backups and immutable archives.

CryptoShred Health demonstrates that cryptographic erasure resolves this deadlock by decoupling data persistence from cryptographic accessibility. Converting erasure into a verifiable cryptographic operation guarantees irreversible data destruction across distributed storage tiers in invariant constant time. Through post-quantum digital certificates, dynamic retention horizons, and automated disaster recovery reconciliation, zero-plaintext architectures demonstrate that healthcare systems can enforce verifiable privacy guarantees and information-theoretic irrecoverability while sustaining high-throughput clinical performance and long-term regulatory auditability.